% Generated draft: requires mathematical and source review.
\documentclass[11pt]{article}
\usepackage[margin=25mm]{geometry}
\usepackage{fontspec}
\setmainfont{Latin Modern Roman}
\usepackage{amsmath,amssymb,mathtools,unicode-math}
\setmathfont{Latin Modern Math}
\usepackage{enumitem,needspace,xurl,hyperref}
\setlength{\parindent}{0pt}
\setlength{\parskip}{5pt}
\setlength{\emergencystretch}{3em}
\newcommand{\N}{\mathbb{N}}
\newcommand{\Z}{\mathbb{Z}}
\newcommand{\Q}{\mathbb{Q}}
\newcommand{\R}{\mathbb{R}}
\newcommand{\C}{\mathbb{C}}
\newcommand{\F}{\mathbb{F}}
\newcommand{\A}{\mathbb{A}}
\newcommand{\PP}{\mathbb{P}}
\newcommand{\E}{\mathbb{E}}
\newcommand{\eps}{\varepsilon}
\newcommand{\dd}{\,\mathrm{d}}
\begin{document}
\section*{UnsolvedMath ID: 30004272}
\textbf{Source problem code:} catalogue entry
\subsection{Definitions and mathematical statement}
The supplied record does not specify the domain of the parameter \(\ell\). The following gives a precise formulation for the positive-integer interpretation; whether this is the full intended parameter range requires source review.

For each \(\ell\in\Z_{\geq 1}\), set
\[
a_\ell(x,z)=x(x^{2\ell}+z^2),\qquad X_1=\partial_x,\qquad X_2=a_\ell(x,z)\partial_z.
\]
The regular region, where these vector fields are linearly independent, is
\[
M_\ell=\{(x,z)\in\R^2:a_\ell(x,z)\ne0\}=\{(x,z)\in\R^2:x\ne0\}.
\]
Interpret the associated Riemannian metric as the metric making \(X_1,X_2\) an orthonormal frame on \(M_\ell\). Thus
\[
g_\ell=\dd x^2+\frac{\dd z^2}{a_\ell(x,z)^2},\qquad
\dd\mu_\ell=\frac{\dd x\,\dd z}{|a_\ell(x,z)|}.
\]
Use the Laplace--Beltrami convention
\[
\Delta_{g_\ell}f
=|a_\ell|\,\partial_x\!\left(|a_\ell|^{-1}\partial_x f\right)
 +|a_\ell|\,\partial_z\!\left(|a_\ell|\partial_z f\right).
\]
Let \(H_\ell=-\Delta_{g_\ell}\) be the densely defined symmetric operator in the complex Hilbert space \(L^2(M_\ell,\mu_\ell)\), with initial domain
\[
\mathcal D(H_\ell)=C_c^\infty(M_\ell;\C).
\]
Here the entire regular region, including both components \(x>0\) and \(x<0\), is used. No boundary condition at the omitted line \(x=0\), or at infinity, is imposed on this initial operator. Essential self-adjointness means that its closure is self-adjoint, equivalently \(\overline{H_\ell}=H_\ell^*\).

\textbf{Question, under the positive-integer interpretation.} Is \(H_\ell\) essentially self-adjoint for every \(\ell\in\Z_{\geq1}\)?

The metric, volume measure, and initial operator domain above make explicit the conventional interpretation of the supplied Laplace--Beltrami question; the record itself does not state these choices explicitly. Its notation \(\ell\geq1\) does not specify whether noninteger parameters are intended. Such an extension requires a definition of \(x^{2\ell}\) for negative \(x\) and appropriate regularity assumptions, which are not supplied. No convention such as replacing \(x^{2\ell}\) by \(|x|^{2\ell}\) is adopted here.

The supplied dated assessment labels the problem partially solved and reports that the case \(\ell=1\) is proved, whereas the corresponding assertion for all positive \(\ell\) was expected but not proved in the report. These are supplied literature claims, not independently verified conclusions. The original record also mentioned a preceding real-analytic no-tangency conjecture; the supplied statement-verification note says that this broader question duplicated \(\texttt{OWR-17290-001}\) and was removed. It is therefore not included as a second target here.
\subsection{Short English statement}
For each positive integer \(\ell\), give the region \(x\ne0\) the Riemannian metric for which \(\partial_x\) and \(x(x^{2\ell}+z^2)\partial_z\) are orthonormal. Ask whether its Laplace--Beltrami operator on compactly supported smooth functions is essentially self-adjoint. The intended parameter domain needs verification. Supplied dated curation reports a proof for \(\ell=1\) and labels the general question partially solved; this assessment has not been independently verified.
\subsection{Sources}
\begin{itemize}
\item Local UnsolvedMath registry, supplied record with permanent ID \(30004272\), title \emph{Self-Adjointness for Degenerate No-Tangency Metrics}, source problem code \texttt{catalogue entry}, in \emph{Oberwolfach Reports Open Problems}. The record supplies proposed year 2019 and no named proposer. Its statement-status label is \texttt{corrected\_verified}; that label is supplied registry metadata, not independent verification of this draft.
\item Supplied report citation: \emph{Mini-Workshop: Self-adjoint Extensions in New Settings} (2020), \emph{Oberwolfach Reports}, pp. 2911--2950, DOI \texttt{10.4171/owr/2019/47}. \url{https://doi.org/10.4171/owr/2019/47}. The record describes these reports as documenting mathematical workshops and their open-problem discussions.
\item Supplied literature source, labeled ``OWR 2019 47'': \url{https://ems.press/content/serial-article-files/46828}.
\item Supplied literature source, attributed to Franceschi--Prandi--Rizzi: \url{https://doi.org/10.1007/s11118-018-09760-w}. No article title, theorem number, or precise supporting passage is supplied.
\item Literature assessment date: \texttt{2026-08-22}. This is supplied dated curation, not independently verified. It reports that the explicit family is proved for \(\ell=1\), while the report expected but did not prove the analogous essential-self-adjointness result for all positive \(\ell\). The supplied background expressly requests independent verification of time-sensitive status and cited claims.
\item Supplied statement-verification note: the general no-tangency question duplicated \texttt{OWR-17290-001} and was removed; the distinct degenerate model was retained.
\end{itemize}
\end{document}
